% Part: normal-modal-logic
% Chapter: completeness
% Section: frame-completeness

\documentclass[../../../include/open-logic-section]{subfiles}

\begin{document}

\olfileid{nml}{com}{fra}

\olsection{चौकट-संपूर्णता}

दिलेल्या तर्कशास्त्राच्या कॅनॉनिकल प्रतिमानाला
संबंधित चौकट-गुणधर्म आहे हे दाखवल्यावर
\Log{K} साठीचे संपूर्णता प्रमेय इतर
मोडल प्रणालींनाही लागू करता येते.

\begin{thm}\ollabel{thm:completeframeprops}
  सामान्य मोडल तर्कशास्त्र~$\Sigma$ मध्ये
  \olref{tab:correspondencetable} च्या
  डाव्या बाजूच्या !!{formula}s पैकी एखादे असेल,
  तर~$\Sigma$ साठीच्या कॅनॉनिकल प्रतिमानाला
  उजव्या बाजूचा संबंधित गुणधर्म असतो.
\end{thm}

\begin{table}[htp]
  \centering
    \begin{tabular}{| l || l |}
      \hline
      \emph{$\Sigma$ मध्ये \dots असेल, तर} & \emph{$\Sigma$ साठीचे
        कॅनॉनिकल प्रतिमान \dots असते:} \\
      \hline \hline
      \Ax{D}:\;\; $\Box!A \lif \Diamond !A$ & \quad सर्वत्र उत्तरवर्ती; \\
      \hline
      \Ax{T}:\;\; $\Box!A \lif !A$ &\quad  परावर्ती;\\
      \hline
      \Ax{B}: \;\;$!A \lif \Box\Diamond!A$ &\quad  सममित; \\
      \hline
      \Ax{4}: \;\; $\Box!A \lif \Box\Box!A$ & \quad संक्रमक; \\
      \hline
      \Ax{5}: \;\; $\Diamond !A \lif \Box\Diamond!A$& \quad युक्लिडी.\\
      \hline
    \end{tabular}
    \caption{मूलभूत अनुरूपता तथ्ये.}
    \ollabel{tab:correspondencetable}
\end{table}

\begin{proof}
  यांपैकी प्रत्येक प्रकरण स्वतंत्रपणे पाहू.

  $\Sigma$ मध्ये \Ax{D} आहे असे समजा आणि
  $\Delta \in W^\Sigma$ घ्या.
  $R^\Sigma \Delta\Delta'$ असा
  $\Delta'$ अस्तित्वात आहे हे दाखवायचे आहे.
  त्यासाठी $\Box^{-1}\Delta$ हा
  $\Sigma$-सुसंगत आहे हे दाखवणे पुरेसे;
  कारण मग लिंडेनबाउमच्या पूर्वप्रमेयानुसार
  त्याचा विस्तार असलेला संपूर्ण
  $\Sigma$-सुसंगत संच
  $\Delta' \supseteq \Box^{-1}\Delta$
  असतो आणि $R^\Sigma$ च्या व्याख्येनुसार
  \iftag{prvBox}{}{आणि \olref[mod]{lem:box-iff-diamond} नुसार }
  $R^\Sigma \Delta\Delta'$.
  म्हणून व्याघातासाठी $\Box^{-1}\Delta$
  हा $\Sigma$-सुसंगत \emph{नाही}
  असे समजा; म्हणजे
  $\Box^{-1}\Delta \Proves[\Sigma] \lfalse$.
  \olref[mod]{lem:box2} नुसार
  $\Delta \Proves[\Sigma] \Box\lfalse$;
  आणि $\Sigma$ मध्ये \Ax{D} असल्याने
  $\Delta \Proves[\Sigma] \Diamond\lfalse$
  हेही मिळते. पण $\Sigma$ सामान्य असल्याने
  $\Sigma \Proves \lnot\Diamond \lfalse$
  (\olref[prf][nor]{prop:notDiamondBot});
  म्हणून $\Delta
  \Proves[\Sigma] \lnot\Diamond \lfalse$
  हेही मिळते. हा~$\Delta$ च्या
  सुसंगततेशी व्याघात आहे.

  आता $\Sigma$ मध्ये \Ax{T} आहे
  असे समजा आणि $\Delta \in W^\Sigma$
  घ्या. $R^\Sigma \Delta\Delta$
  दाखवायचे आहे; म्हणजेच
  \iftag{prvBox}{}{\olref[mod]{lem:box-iff-diamond} नुसार }
  $\Box^{-1}\Delta \subseteq \Delta$.
  पण $\Box!A \in \Delta$ असेल,
  तर \Ax{T} नुसार $!A \in \Delta$;
  हेच अपेक्षित होते.

  आता $\Sigma$ मध्ये \Ax{B} आहे
  असे समजा आणि $\Delta$, $\Delta' \in
  W^\Sigma$ साठी $R^\Sigma
  \Delta\Delta'$ असे धरा.
  $R^\Sigma \Delta'\Delta$
  दाखवायचे आहे; म्हणजे
  \iftag{prvBox}{$\Box^{-1}\Delta'
    \subseteq \Delta$.
    \olref[mod]{lem:box-iff-diamond} नुसार
    ही अट पुढील अटीशी समतुल्य आहे:}{}
  $\Diamond\Delta \subseteq \Delta'$.
  म्हणून $!A \in \Delta$ असे
  समजा. \Ax{B} नुसार
  $\Box\Diamond!A \in \Delta$.
  $R^\Sigma \Delta\Delta'$ या
  गृहीतकावरून \iftag{prvBox}{}{आणि
    \olref[mod]{lem:box-iff-diamond} नुसार}
  $\Box^{-1}\Delta \subseteq \Delta'$;
  म्हणून आवश्यकतेनुसार
  $\Diamond!A \in \Delta'$.

  आता $\Sigma$ मध्ये \Ax{4} आहे असे
  समजा आणि $R^\Sigma \Delta_1\Delta_2$
  व $R^\Sigma \Delta_2\Delta_3$
  असे धरा. $R^\Sigma
  \Delta_1\Delta_3$ दाखवायचे आहे.
  गृहीतकावरून \iftag{prvBox}{}{आणि
    \olref[mod]{lem:box-iff-diamond} नुसार}
  $\Box^{-1}\Delta_1 \subseteq \Delta_2$
  आणि $\Box^{-1}\Delta_2 \subseteq \Delta_3$
  या दोन्ही अटी मिळतात.
  $R^\Sigma \Delta_1\Delta_3$
  साठी $\Box^{-1}\Delta_1
  \subseteq \Delta_3$ दाखवणे पुरेसे.
  म्हणून $!B \in \Box^{-1}\Delta_1$
  घ्या; म्हणजे $\Box!B \in \Delta_1$.
  \Ax{4} नुसार $\Box\Box!B \in
  \Delta_1$ हेही मिळते. गृहीतकांवरून
  आधी $\Box!B \in \Delta_2$
  आणि नंतर $!B \in \Delta_3$,
  हेच अपेक्षित होते.

  आता $\Sigma$ मध्ये \Ax{5} आहे
  असे समजा आणि $R^\Sigma
  \Delta_1\Delta_2$ व
  $R^\Sigma \Delta_1\Delta_3$
  असे धरा. $R^\Sigma
  \Delta_2\Delta_3$ दाखवायचे आहे.
  पहिल्या गृहीतकावरून
  $\Box^{-1}\Delta_1 \subseteq \Delta_2$
  \iftag{prvBox}{}{\olref[mod]{lem:box-iff-diamond} नुसार}
  मिळते; दुसरे गृहीतक
  $\Diamond\Delta_3 \subseteq \Delta_1$
  याच्याशी समतुल्य आहे
  \iftag{prvBox}{\olref[mod]{lem:box-iff-diamond} नुसार}{}.
  $R^\Sigma \Delta_2\Delta_3$
  दाखवण्यासाठी \iftag{prvBox}{
    \olref[mod]{lem:box-iff-diamond} नुसार,
    पुढील अट दाखवणे पुरेसे:}{पुढील अट दाखवायची आहे:}
  $\Diamond\Delta_3 \subseteq \Delta_2$.
  म्हणून $\Diamond!A \in
  \Diamond\Delta_3$ घ्या;
  म्हणजे $!A \in \Delta_3$.
  दुसऱ्या गृहीतकानुसार
  $\Diamond!A \in \Delta_1$;
  \Ax{5} नुसार
  $\Box\Diamond!A \in \Delta_1$
  हेही मिळते. आता पहिल्या
  गृहीतकावरून $\Diamond!A \in
  \Delta_2$, हेच अपेक्षित होते.
  \end{proof}

या निष्कर्षांवरून अनेक प्रणालींसाठी
संपूर्णतेची निष्पन्नता मिळते.
उदाहरणार्थ, $\Log{S5} = \Log{KT5} =
\Log{KTB4}$ हे सर्व परावर्ती युक्लिडी
प्रतिमानांच्या वर्गाच्या सापेक्ष संपूर्ण
आहे, हे आपल्याला माहीत आहे.
हा वर्ग सर्व परावर्ती, सममित आणि
संक्रमक प्रतिमानांच्या वर्गाइतकाच आहे.

\begin{thm}\ollabel{thm:generaldet}
  $\mClass{C}_\Ax{D}$,
  $\mClass{C}_\Ax{T}$,
  $\mClass{C}_\Ax{B}$,
  $\mClass{C}_\Ax{4}$ आणि
  $\mClass{C}_\Ax{5}$
  हे अनुक्रमे सर्व सर्वत्र उत्तरवर्ती,
  परावर्ती, सममित, संक्रमक आणि
  युक्लिडी प्रतिमानांचे वर्ग असू द्या.
  मग \Ax{D}, \Ax{T}, \Ax{B},
  \Ax{4} आणि \Ax{5} यांमधील
  कोणत्याही योजना $!A_1$, \dots, $!A_n$
  साठी प्रणाली
  $\Log{K}!A_1 \dots !A_n$
  हिचे निर्धारण प्रतिमानांचा वर्ग
  $\mClass{C} = \mClass{C}_{!A_1} \cap \dots
  \cap \mClass{C}_{!A_n}$ करतो.
\end{thm}

\begin{prop}
  $\Sigma$ हे सामान्य मोडल तर्कशास्त्र
  असू द्या. मग:
  \begin{enumerate}
  \item\ollabel{prop:anotherfive-a}%
    $\Sigma$ मध्ये $\Diamond!A \lif \Box
    !A$ ही योजना असेल, तर~$\Sigma$
    साठीचे कॅनॉनिकल प्रतिमान
    आंशिक फलनीय असते.
  \item $\Sigma$ मध्ये
    $\Diamond!A \liff \Box
    !A$ ही योजना असेल, तर~$\Sigma$
    साठीचे कॅनॉनिकल प्रतिमान
    फलनीय असते.
  \item $\Sigma$ मध्ये
    $\Box\Box!A \lif \Box
    !A$ ही योजना असेल, तर~$\Sigma$
    साठीचे कॅनॉनिकल प्रतिमान
    दुर्बल सघन असते.
  \end{enumerate}
  (या चौकट-गुणधर्मांच्या व्याख्यांसाठी
  \olref[frd][acc]{tab:anotherfive} पाहा.)
\end{prop}

\begin{proof}
  \begin{enumerate}
  \item $\Sigma$ मध्ये
    $\Diamond !A \lif \Box !A$
    ही योजना आहे असे समजा.
    $R^\Sigma$ आंशिक फलनीय आहे हे
    दाखवण्यासाठी कोणत्याही
    $\Delta_1$, $\Delta_2$,
    $\Delta_3 \in W^\Sigma$
    साठी $R^\Sigma \Delta_1\Delta_2$
    आणि $R^\Sigma \Delta_1\Delta_3$
    असतील, तर $\Delta_2=\Delta_3$
    हे सिद्ध करायचे आहे.
    $R^\Sigma \Delta_1\Delta_2$
    असल्याने $\Box^{-1}\Delta_1
    \subseteq \Delta_2$;
    आणि $R^\Sigma \Delta_1\Delta_3$
    असल्याने $\Box^{-1}\Delta_1
    \subseteq \Delta_3$
    \iftag{prvBox}{}{हे दोन्ही
      \olref[mod]{lem:box-iff-diamond} नुसार}.
    $\Delta_2=\Delta_3$ साठी
    $\Delta_2 \subseteq \Delta_3$
    आणि $\Delta_3 \subseteq \Delta_2$
    हे दोन्ही समावेश सिद्ध करणे पुरेसे.
    पहिल्या समावेशासाठी
    $!A \in \Delta_2$ घ्या;
    मग $\Diamond!A \in \Delta_1$.
    या योजनेवरून आणि~$\Delta_1$ च्या
    निगामी बंदतेवरून
    $\Box!A \in \Delta_1$ हेही मिळते.
    $R^\Sigma \Delta_1\Delta_3$
    या गृहीतकावरून $!A \in \Delta_3$.
    दुसरा समावेशही तसाच सिद्ध होतो.

  \item हे \olref{prop:anotherfive-a}
    आणि \olref{thm:completeframeprops}
    मधील सर्वत्र उत्तरवर्तित्वाच्या
    सिद्धतेवरून लगेच मिळते.

  \item $\Sigma$ मध्ये
    $\Box\Box!A \lif \Box!A$
    ही योजना आहे असे समजा.
    $R^\Sigma$ दुर्बल सघन आहे हे
    दाखवण्यासाठी $R^\Sigma
    \Delta_1\Delta_2$ घ्या.
    $R^\Sigma \Delta_1\Delta_3$
    आणि $R^\Sigma \Delta_3\Delta_2$
    असा संपूर्ण $\Sigma$-सुसंगत
    संच~$\Delta_3$ अस्तित्वात आहे
    हे दाखवायचे आहे. असे ठेवू:
    \[
    \Gamma = \Box^{-1}\Delta_1 \cup \Diamond\Delta_2.
    \]
    $\Gamma$ हा $\Sigma$-सुसंगत
    आहे हे दाखवणे पुरेसे. मग
    लिंडेनबाउमच्या पूर्वप्रमेयानुसार
    त्याचा विस्तार असा संपूर्ण
    $\Sigma$-सुसंगत संच~$\Delta_3$
    मिळतो की $\Box^{-1}\Delta_1
    \subseteq \Delta_3$ आणि
    $\Diamond\Delta_2 \subseteq \Delta_3$.
    म्हणजे $R^\Sigma \Delta_1\Delta_3$
    \iftag{prvBox}{}{(
      \olref[mod]{lem:box-iff-diamond} नुसार)}
    आणि $R^\Sigma \Delta_3\Delta_2$
    \iftag{prvBox}{(
      \olref[mod]{lem:box-iff-diamond} नुसार)}{}.

    व्याघातासाठी $\Gamma$ सुसंगत
    नाही असे समजा. मग
    $\Box!A_1$, \dots,
    $\Box!A_n \in \Delta_1$
    आणि $!B_1$, \dots,~$!B_m \in \Delta_2$
    अशी !!{formula}s असतात की
    \[!A_1, \dots,
    !A_n, \Diamond!B_1, \dots, \Diamond!B_m \Proves[\Sigma]
    \lfalse.\]
    $\Diamond (!B_1 \land \dots \land !B_m) \to
    (\Diamond!B_1 \land \dots \land \Diamond!B_m)$
    हे प्रत्येक सामान्य मोडल तर्कशास्त्रात
    !!{derivable} असल्याने
    पुढीलप्रमाणे युक्तिवाद करून
    आपण~$\Delta_2$ च्या
    सुसंगततेशी व्याघात मिळवतो:
    \begin{align*}
      !A_1, \dots, !A_n, & \Diamond!B_1, \dots,
      \Diamond!B_m \Proves[\Sigma] \lfalse \\
      !A_1, \dots,!A_n & \Proves[\Sigma] (\Diamond!B_1 \land
      \dots \land \Diamond!B_m) \lif \lfalse\\
      & \qquad\text{निगमन प्रमेयानुसार}\\
      & \qquad\text{\olref[prf][prp]{prop:derivabilityfacts}\olref[prf][prp]{prop:derivabilityfacts-deduction}, आणि \Taut} \\
      !A_1, \dots,!A_n & \Proves[\Sigma]
      \Diamond(!B_1 \land
      \dots \land !B_m) \lif \lfalse\\
      & \qquad \text{$\Sigma$ सामान्य असल्याने} \\
      !A_1, \dots,!A_n & \Proves[\Sigma]
      \lnot\Diamond (!B_1 \land \dots \land !B_m)\\
      & \qquad\text{\PL नुसार} \\
      !A_1, \dots,!A_n & \Proves[\Sigma]
      \Box\lnot (!B_1 \land \dots \land !B_m)\\
      & \qquad\text{$\lnot\Diamond$ साठी $\Box\lnot$} \\
      \Box!A_1, \dots,\Box!A_n
      & \Proves[\Sigma] \Box\Box \lnot (!B_1 \land \dots \land !B_m)\\
      & \qquad\text{\olref[mod]{lem:box1} नुसार} \\
      \Box!A_1, \dots,\Box!A_n
      & \Proves[\Sigma]
      \Box\lnot (!B_1 \land \dots \land !B_m)\\
      &\qquad\text{$\Box\Box!A \lif \Box!A$ योजनेनुसार} \\
      \Delta_1 & \Proves[\Sigma]
      \Box\lnot (!B_1 \land \dots \land !B_m)\\
      &\qquad\text{एकस्वनिकतेनुसार, \olref[prf][prp]{prop:derivabilityfacts}\olref[prf][prp]{prop:derivabilityfacts-monotonicity}} \\
      & \Box\lnot (!B_1 \land \dots \land !B_m) \in \Delta_1\\
      &\qquad\text{निगामी बंदतेनुसार}; \\
      & \lnot (!B_1 \land \dots \land !B_m) \in \Delta_2\\
      &\qquad \text{कारण }
      R^\Sigma \Delta_1\Delta_2.
    \end{align*}
  \end{enumerate}
\end{proof}

या उदाहरणांवरून प्रत्येक मोडल
प्रणाली~$\Sigma$ \emph{संपूर्ण}
असेल असे वाटू शकते: म्हणजे~$\Sigma$
ची सर्व प्रमेये ज्या प्रत्येक चौकटीत
वैध आहेत, त्या प्रत्येक चौकटीत
वैध असलेले प्रत्येक सूत्रही
ती प्रणाली सिद्ध करत असेल.
पण दुर्दैवाने या अर्थाने संपूर्ण
नसलेल्या अनेक प्रणाली आहेत.

\end{document}
